\chapter*{Lời Mở Đầu}
\addcontentsline{toc}{chapter}{Lời Mở Đầu}
\markboth{Lời Mở Đầu}{Lời Mở Đầu}

\begin{truyen}{Lời Mở Đầu}{Opening Words}
\chuhoa{C}{ó những quyển sách nhìn lớp học như một chỗ để dạy chữ. Quyển này nhìn lớp học như một nơi va chạm: va giữa tuổi trẻ với kỳ vọng, giữa lý tưởng của người lớn với nhịp sống của Gen Z, giữa tình thương của cha mẹ với áp lực thành tích, và giữa sự nghiêm khắc của nhà trường với cái tôi đang lớn của học sinh.}
\textit{(Some books look at the classroom as a place to teach knowledge. This book looks at it as a place of collision: between youth and expectations, adult ideals and Gen Z rhythms, parental love and achievement pressure, and school discipline and the growing ego of students.)}

\textbf{Đa Giác Học Đường} đi theo bốn phần: \textbf{Tâm lý chiến Gen Z}, \textbf{Lý sự xoay chiều}, \textbf{Cuộc chiến ba tròn}, và \textbf{Ảo tưởng tương lai}. Bốn phần ấy không cố phân xử ai đúng tuyệt đối. Chúng chỉ bóc từng lớp ngụy biện để người đọc thấy rằng mỗi câu cãi trong lớp thường đứng trên một nỗi sợ có thật.
\textit{(\textbf{The Classroom Multiverse} moves through four parts: \textbf{Gen Z Mental Warfare}, \textbf{Turning the Tables}, \textbf{The Three-Way War}, and \textbf{Future Delusions}. These four parts do not try to declare one side absolutely right. They peel away layers of excuses to show that many classroom arguments stand on top of a real fear.)}

Nếu đọc xong mà học sinh bớt đổ lỗi, phụ huynh bớt nóng, và thầy cô bớt nghĩ rằng chỉ cần đúng lý là đủ, thì quyển này đã hoàn thành việc của nó.
\textit{(If, after reading, students blame less, parents react less harshly, and teachers realize that logic alone is not enough, then this book has done its work.)}
\end{truyen}

\begin{baihoc}
Muốn hiểu một câu nói trong lớp, đừng chỉ nghe chữ. Hãy nghe xem đằng sau nó là sĩ diện, kiệt sức, áp lực hay sự thiếu trưởng thành.
\textit{(To understand a sentence in class, do not only hear the words. Hear the pride, exhaustion, pressure, or immaturity behind it.)}
\end{baihoc}
\ngancach